\section{Phương pháp đề xuất}
\label{sec:method}

\subsection{Tổng quan pipeline}

Pipeline được thiết kế thành các bước độc lập nhưng trao đổi với nhau thông
qua các artifact JSON có schema xác định. Cách tổ chức này cho phép kiểm tra
từng giai đoạn, chạy lại một bước mà không phải thực thi toàn bộ pipeline và
so sánh các phương pháp xác minh trên cùng kết quả retrieval.

Luồng xử lý tổng thể gồm năm giai đoạn:

\begin{enumerate}
    \item chuẩn hóa văn bản pháp luật thành các quy tắc điều
    kiện--hệ quả;

    \item xây dựng đồ thị nhân quả pháp luật và causal memory;

    \item truy hồi event, rule và causal path nhiều bước;

    \item phân tích claim theo truy vấn và xác minh bằng LegalSCM hoặc path
    ablation;

    \item lựa chọn evidence, xác định answer intent và sinh câu trả lời có
    citation.
\end{enumerate}

Với truy vấn đầu vào \(q\), pipeline tạo ra một tập candidate event, một tập
candidate rule và một tập causal path. Các path sau đó được xác minh để chọn
primary path. Câu trả lời cuối cùng được sinh từ primary path, quyết định
xác minh và những evidence pháp luật đã được giữ lại.

\begin{figure*}[t]
    \centering
    % Thay bằng hình pipeline chính thức khi hoàn thiện paper:
    % \includegraphics[width=0.95\textwidth]{figures/causalrag_pipeline.pdf}
    \fbox{
        \parbox{0.92\textwidth}{
        \centering
        Văn bản pháp luật
        \(\rightarrow\) Chuẩn hóa quy tắc
        \(\rightarrow\) Đồ thị nhân quả và causal memory
        \(\rightarrow\) Truy hồi đa bước
        \(\rightarrow\) Query-aware verification
        \(\rightarrow\) Sinh câu trả lời và citation
        }
    }
    \caption{Kiến trúc tổng thể của pipeline CausalRAG cho suy luận pháp
    luật hình sự Việt Nam.}
    \label{fig:pipeline}
\end{figure*}

\subsection{Chuẩn hóa quy tắc pháp luật}

Đầu vào ban đầu là các điều khoản pháp luật ở dạng văn bản. Mục tiêu của
bước chuẩn hóa là chuyển nội dung văn bản thành các đơn vị có thể được sử
dụng trong retrieval và graph reasoning. Mỗi quy tắc được lưu với các trường
chính:

\begin{itemize}
    \item mã định danh quy tắc;
    \item mã và tiêu đề điều luật;
    \item chủ thể pháp luật;
    \item điều kiện áp dụng;
    \item hệ quả pháp lý;
    \item event ID và event name của điều kiện;
    \item event ID và event name của hệ quả;
    \item loại quan hệ nhân quả hoặc loại hệ quả nếu có.
\end{itemize}

Event ID được chuẩn hóa thành chuỗi không dấu và dùng dấu gạch dưới để kết
nối các token. Event name giữ lại biểu diễn tiếng Việt có thể đọc được. Việc
tách ID khỏi name giúp hệ thống sử dụng ID ổn định cho graph matching, đồng
thời sử dụng name cho semantic retrieval và sinh câu trả lời.

Ví dụ, một quy tắc có thể được chuẩn hóa thành:

\begin{quote}
Điều kiện: ``cố ý thực hiện tội phạm nhưng không thực hiện được đến cùng vì
nguyên nhân ngoài ý muốn'';

Hệ quả: ``được xác định là phạm tội chưa đạt'';

Căn cứ: Điều 15 Bộ luật Hình sự.
\end{quote}

Quy tắc này tạo quan hệ:

\[
\texttt{co\_y\_thuc\_hien\_toi\_pham\_nhung\_khong\_hoan\_thanh}
\rightarrow
\texttt{pham\_toi\_chua\_dat}.
\]

Nếu một event vừa là hệ quả của quy tắc trước vừa là điều kiện của quy tắc
sau, hai rule có thể được nối thành một causal path nhiều bước.

\subsection{Xây dựng đồ thị nhân quả dị thể}

Từ tập rule đã chuẩn hóa, hệ thống tạo ba nhóm node chính: event, rule và
article. Mỗi rule được kết nối với event điều kiện, event hệ quả và điều luật
nguồn. Đồng thời, một cạnh \texttt{CAUSES} được tạo từ event điều kiện đến
event hệ quả.

Việc lưu rule dưới dạng node riêng biệt giúp bảo toàn thông tin pháp lý mà
một event-to-event edge đơn giản có thể làm mất, chẳng hạn:

\begin{itemize}
    \item nhiều quy tắc cùng nối một cặp event;
    \item các quy tắc có cùng outcome nhưng khác điều kiện;
    \item cùng một event xuất hiện trong nhiều điều luật;
    \item một quan hệ được hỗ trợ bởi nhiều căn cứ pháp luật.
\end{itemize}

Đối với graph traversal, hệ thống sử dụng projection trên các event và cạnh
\texttt{CAUSES}. Tuy nhiên, mỗi cạnh vẫn lưu rule ID và article ID để có thể
truy ngược từ path về evidence gốc.

Từ graph đầy đủ, hệ thống đồng thời trích xuất các chuỗi hai bước:

\[
v_0 \rightarrow v_1 \rightarrow v_2.
\]

Các chuỗi này được dùng để kiểm tra cấu trúc graph, hỗ trợ xây dựng benchmark
và tạo candidate path trong retrieval.

\subsection{Xây dựng causal memory}

Graph traversal yêu cầu một seed event phù hợp với truy vấn. Tuy nhiên,
event name trong graph không nhất thiết xuất hiện nguyên văn trong câu hỏi.
Do đó, hệ thống xây dựng một causal memory chứa cả rule record và event
record.

Mỗi rule record bao gồm:

\begin{itemize}
    \item nội dung điều kiện và hệ quả;
    \item chủ thể;
    \item điều luật;
    \item condition event và effect event;
    \item metadata liên kết với graph.
\end{itemize}

Mỗi event record bao gồm:

\begin{itemize}
    \item event ID và event name;
    \item vai trò điều kiện, mediator hoặc hệ quả;
    \item các rule và article liên quan;
    \item graph node ID tương ứng.
\end{itemize}

Các record được chuyển thành biểu diễn vector bằng BGE-M3
\cite{chen2024bgem3}. Chỉ mục vector được xây dựng bằng FAISS để hỗ trợ tìm
kiếm lân cận gần trên tập memory
\cite{johnson2019billion}. Trong thực nghiệm, cùng một embedding space được
sử dụng cho truy vấn, rule text và event text.

Với truy vấn \(q\), embedding truy vấn được tính bởi:

\begin{equation}
\mathbf{z}_q
=
f_{\theta}(q),
\end{equation}

trong đó \(f_{\theta}\) là encoder BGE-M3. Độ tương đồng ngữ nghĩa giữa truy
vấn và memory record \(m_i\) được tính bằng cosine similarity:

\begin{equation}
s_{\mathrm{sem}}
\left(
q,m_i
\right)
=
\frac{
    \mathbf{z}_q^\top \mathbf{z}_{m_i}
}{
    \left\Vert\mathbf{z}_q\right\Vert
    \left\Vert\mathbf{z}_{m_i}\right\Vert
}.
\end{equation}

Semantic score chỉ được sử dụng để tạo và xếp hạng ứng viên. Việc một record
có semantic score cao không tự động làm cho nó trở thành evidence cuối cùng.

\subsection{Truy hồi event và rule ứng viên}

Bước retrieval thực hiện hai nhánh song song.

Nhánh thứ nhất truy hồi các event có độ tương đồng cao với truy vấn. Các
event này được dùng làm seed cho graph traversal. Mỗi seed record lưu:

\begin{itemize}
    \item event ID;
    \item graph node ID;
    \item event name;
    \item semantic score;
    \item vai trò của event;
    \item phương pháp retrieval.
\end{itemize}

Nhánh thứ hai truy hồi trực tiếp các rule có nội dung gần với truy vấn. Direct
rule retrieval giúp giữ lại những rule phù hợp ngay cả khi event name không
khớp tốt với cách diễn đạt trong câu hỏi.

Các ứng viên có điểm dưới ngưỡng tối thiểu bị loại bỏ. Trong cấu hình chính,
hệ thống lấy tối đa tám seed event và tám direct rule hit trước khi mở rộng
graph. Một semantic pool lớn hơn được giữ tạm thời để tránh loại bỏ sớm các
ứng viên có thể tham gia vào một path tốt.

\subsection{Mở rộng causal path nhiều bước}

Từ mỗi seed event, retriever duyệt các cạnh \texttt{CAUSES} theo hướng được
cấu hình. Thực nghiệm sử dụng \texttt{direction=both}, cho phép tìm kiếm theo
chiều thuận và chiều ngược. Tùy theo vị trí của event trong truy vấn, một
seed có thể đóng vai trò nguyên nhân, mediator hoặc outcome.

Quá trình mở rộng bị giới hạn bởi:

\begin{itemize}
    \item số hop tối đa;
    \item số path tối đa trên mỗi seed event;
    \item số candidate rule tối đa;
    \item ngưỡng semantic score của event và rule.
\end{itemize}

Mỗi candidate path lưu:

\begin{itemize}
    \item danh sách event ID và event name theo thứ tự nhân quả;
    \item danh sách rule ID hỗ trợ từng hop;
    \item danh sách article ID;
    \item seed event và outcome event;
    \item số hop;
    \item semantic và graph score;
    \item giải thích cấu trúc của path.
\end{itemize}

Retriever ưu tiên path có cấu trúc liên tục, rule đầy đủ cho từng hop và mức
liên quan cao với truy vấn. Các path trùng nhau được hợp nhất theo chuỗi event
và rule. Kết quả retrieval cuối cùng gồm candidate path, evidence rule và
các event đã truy hồi.

Cách thiết kế này tách biệt hai khái niệm:

\begin{itemize}
    \item \emph{path retrieval}: gold path có xuất hiện trong tập ứng viên
    hay không;

    \item \emph{path ranking}: gold path có được xếp ở vị trí đầu tiên hay
    không.
\end{itemize}

Sự phân biệt này được phản ánh trong hai metric Oracle Exact Path và Top-1
Exact Path ở Phần~\ref{sec:results}.

\subsection{Phân tích truy vấn theo claim}

Kết quả retrieval được chuyển sang bước query-aware verification. Trước khi
xác minh path, hệ thống chuẩn hóa truy vấn bằng cách chuyển về chữ thường,
loại bỏ khoảng trắng dư và nhận diện các nhóm tín hiệu ngôn ngữ.

Các nhóm tín hiệu chính gồm:

\begin{itemize}
    \item tín hiệu về quan hệ trực tiếp, chẳng hạn ``trực tiếp'' hoặc
    ``cạnh trực tiếp'';

    \item tín hiệu loại bỏ hoặc giả định vắng mặt của mediator;

    \item tín hiệu outcome vẫn giữ nguyên;

    \item tín hiệu outcome biến mất;

    \item tín hiệu hỏi về tính cần thiết;

    \item tín hiệu phản thực tế hoặc điều kiện giả định.
\end{itemize}

Từ các tín hiệu này, query analyzer tạo một object
\texttt{query\_analysis} chứa:

\begin{itemize}
    \item normalized query;
    \item claim type;
    \item phiên bản intent detector;
    \item các cờ direct, remove, remains, disappears và necessary;
    \item mediator được ánh xạ;
    \item factual outcome và counterfactual outcome nếu có;
    \item phương pháp xác minh;
    \item trạng thái structural fallback.
\end{itemize}

Đối với câu hỏi phản thực tế, mediator được ánh xạ tới event trên candidate
path bằng exact matching hoặc semantic matching. Nếu không tìm được mediator
đủ tin cậy, hệ thống không tự động giả định một event tùy ý mà giữ kết quả ở
trạng thái chưa xác định hoặc sử dụng fallback có ghi nhận lý do.

\subsection{Khởi tạo LegalSCM}

LegalSCM được xây dựng trực tiếp từ tập rule đã chuẩn hóa. Mỗi event được ánh
xạ thành một biến nội sinh ba giá trị:

\[
X_v
\in
\left\{
    \mathrm{TRUE},
    \mathrm{FALSE},
    \mathrm{UNKNOWN}
\right\}.
\]

Mỗi rule được chuyển thành một mechanism liên kết condition event với effect
event. Với phiên bản corpus hiện tại, một mechanism cơ bản có dạng:

\begin{equation}
X_{e_i}
\leftarrow
f_{r_i}
\left(
    X_{c_i}
\right).
\end{equation}

Khi \(X_{c_i}=\mathrm{TRUE}\), rule \(r_i\) có thể kích hoạt và suy ra effect
event. Nếu condition chưa xác định, rule chưa đủ căn cứ để kích hoạt.
LegalSCM lặp qua các mechanism cho đến khi đạt fixed point hoặc đạt số vòng
lặp tối đa.

Các event gốc không được gán bởi factual context được giữ ở trạng thái
\(\mathrm{UNKNOWN}\). Đối với event nội sinh, hệ thống áp dụng giả định
closed-world có kiểm soát khi không có mechanism nào có thể kích hoạt nó.
Nếu các mechanism tạo ra kết quả xung đột, trạng thái cuối được đặt thành
\(\mathrm{UNKNOWN}\) thay vì buộc chọn một kết quả.

LegalSCM được cache theo file rule và thời điểm sửa đổi. Cơ chế cache tránh
việc phân tích lại toàn bộ tập 2.884 rule cho từng câu hỏi trong batch.

\subsection{Xác nhận factual causal path}

Với candidate path:

\[
p=
\left(
v_0,r_1,v_1,\ldots,r_h,v_h
\right),
\]

hệ thống đặt seed event \(v_0\) vào factual context và thực hiện suy luận
fixed point. Một path chỉ được xem là có factual support khi:

\begin{enumerate}
    \item outcome event \(v_h\) được suy ra ở trạng thái
    \(\mathrm{TRUE}\);

    \item các rule được kỳ vọng trên path xuất hiện trong tập rule đã kích
    hoạt;

    \item mỗi hop có ít nhất một mechanism phù hợp;

    \item không có xung đột làm cho event quan trọng trở thành
    \(\mathrm{UNKNOWN}\).
\end{enumerate}

Độ bao phủ rule của path được tính bằng:

\begin{equation}
\mathrm{Coverage}(p)
=
\frac{
    \left|
        \mathcal{R}_{p}
        \cap
        \mathcal{R}_{\mathrm{activated}}
    \right|
}{
    \left|
        \mathcal{R}_{p}
    \right|
},
\end{equation}

trong đó \(\mathcal{R}_{p}\) là tập rule được kỳ vọng trên path và
\(\mathcal{R}_{\mathrm{activated}}\) là tập rule được kích hoạt trong factual
world.

Ngoài tỷ lệ bao phủ tổng thể, hệ thống lưu trạng thái bao phủ của từng hop.
Trace này giúp phân biệt trường hợp outcome được suy ra qua đúng path đang
xét với trường hợp outcome được kích hoạt bởi một rule khác ngoài path.

\subsection{Hard intervention trên mediator}

Đối với mỗi mediator \(M\) nằm bên trong một path nhiều bước, LegalSCM xây
dựng một counterfactual world bằng phép:

\[
do
\left(
    M=\mathrm{FALSE}
\right).
\]

Trong world này, mechanism thông thường tạo ra \(M\) bị thay thế bởi phép
gán hằng \(\mathrm{FALSE}\). Những rule có điều kiện phụ thuộc vào \(M\) được
đánh giá lại, sau đó trạng thái của các event hậu duệ được suy luận lại đến
fixed point. Cách thực hiện này tuân theo semantics của intervention trong
SCM \cite{pearl2009causality}.

Với outcome \(Y\), hệ thống so sánh:

\[
Y^{F}
=
Y
\left(
W^{F}
\right)
\]

và:

\[
Y^{CF}
=
Y
\left(
W^{F}
\mid
do(M=\mathrm{FALSE})
\right).
\]

Nếu \(Y^{F}=\mathrm{TRUE}\) và \(Y^{CF}=\mathrm{FALSE}\), mediator được gán
trạng thái \texttt{NECESSARY} trong context đang xét. Nếu outcome vẫn ở
trạng thái \(\mathrm{TRUE}\), mediator không hoàn toàn cần thiết trong mô
hình. Nếu factual hoặc counterfactual outcome không xác định, intervention
được gán trạng thái chưa giải quyết.

Mỗi intervention trace lưu:

\begin{itemize}
    \item mediator ID và mediator name;
    \item intervention assignment;
    \item factual mediator state;
    \item factual outcome;
    \item counterfactual outcome;
    \item outcome có thay đổi hay không;
    \item các rule bị vô hiệu hóa;
    \item các rule mới được kích hoạt;
    \item rule thay thế dẫn tới cùng outcome;
    \item trạng thái và giải thích của intervention.
\end{itemize}

Những trường này được lưu trong output Step 4 nhằm phục vụ audit và không chỉ
dùng để tính một decision label duy nhất.

\subsection{Node-deletion path ablation}

Path ablation được triển khai bằng cách loại mediator khỏi event graph, sau
đó tìm các đường thay thế từ seed event đến outcome. Tìm kiếm đường thay thế
được giới hạn bởi số hop và số path tối đa để kiểm soát chi phí.

Với mỗi alternative path, hệ thống lưu:

\begin{itemize}
    \item chuỗi event thay thế;
    \item số hop;
    \item các rule hỗ trợ;
    \item graph score;
    \item điểm của đường thay thế tốt nhất.
\end{itemize}

Nếu không tìm được đường thay thế, mediator được xem là cần thiết theo tiêu
chí reachability. Nếu một đường thay thế vượt ngưỡng điểm, outcome được xem
là vẫn có thể tiếp cận sau khi xóa mediator.

Khi chạy ở chế độ \texttt{structural\_scm}, kết quả path ablation vẫn được
tính và lưu dưới dạng baseline diagnostic. Nếu LegalSCM không thể khởi tạo
hoặc không thể xác minh path, hệ thống có thể sử dụng path ablation làm
fallback, nhưng phải ghi rõ:

\begin{itemize}
    \item \texttt{structural\_fallback\_used};
    \item lý do fallback;
    \item phương pháp thực sự tạo ra kết quả.
\end{itemize}

Trong paired experiment, chế độ \texttt{path\_ablation} được chạy độc lập
với cùng retrieval output và cùng hyperparameter, giúp so sánh hai engine mà
không thay đổi các thành phần còn lại.

\subsection{Xếp hạng và phân loại evidence}

Mỗi rule evidence nhận một verification score kết hợp bốn tín hiệu:

\begin{itemize}
    \item mức hỗ trợ từ các factual path;
    \item mức hỗ trợ từ counterfactual verification;
    \item semantic relevance với truy vấn;
    \item graph relevance và retrieval score.
\end{itemize}

Evidence thuộc primary causal path nhận thêm điểm ưu tiên. Dựa trên
verification score, evidence được chia thành ba nhóm:

\begin{description}
    \item[\texttt{KEEP}] Evidence đủ mạnh để được chuyển sang bước sinh câu
    trả lời.

    \item[\texttt{UNCERTAIN}] Evidence có liên quan nhưng chưa đủ mức hỗ trợ.

    \item[\texttt{REMOVE}] Evidence không đạt ngưỡng hoặc bị kết quả xác
    minh bác bỏ.
\end{description}

Trong cấu hình thực nghiệm, keep threshold được đặt bằng \(0{,}52\), còn
reject threshold bằng \(0{,}34\). Evidence thuộc primary path được giữ theo
đúng thứ tự rule trên path, thay vì chỉ sắp xếp theo verification score.

\subsection{Chọn primary path và quyết định cuối}

Các path được xếp hạng dựa trên:

\begin{itemize}
    \item trạng thái factual support;
    \item độ đầy đủ của path;
    \item số hop;
    \item consistency score;
    \item graph và retrieval score;
    \item mức phù hợp với loại claim.
\end{itemize}

Hệ thống ưu tiên path ở trạng thái \texttt{SUPPORTED}, có đủ hai hop đối với
câu hỏi nhiều bước và có consistency score cao. Path đứng đầu được chọn làm
primary path, trong khi các path còn lại được giữ làm supporting path hoặc
phục vụ branch/convergence reasoning.

Khác với phiên bản quyết định theo đa số path, phiên bản query-aware tạo
final decision từ loại claim và primary path. Cụ thể:

\begin{itemize}
    \item với claim nhân quả chuẩn, hệ thống kiểm tra factual support của
    primary path;

    \item với direct claim, hệ thống kiểm tra cạnh trực tiếp thay vì chỉ
    kiểm tra reachability;

    \item với mediator claim, hệ thống so sánh factual và counterfactual
    outcome;

    \item nếu mediator mapping hoặc trạng thái world chưa xác định, quyết
    định được chuyển thành \texttt{UNCERTAIN}.
\end{itemize}

Output xác minh gồm:

\begin{itemize}
    \item \texttt{final\_decision};
    \item \texttt{decision\_score};
    \item \texttt{decision\_explanation};
    \item \texttt{primary\_path\_ids};
    \item \texttt{query\_analysis};
    \item \texttt{path\_verifications};
    \item evidence được giữ, chưa chắc chắn và bị loại;
    \item configuration và phiên bản Step 4.
\end{itemize}

\subsection{Suy luận answer intent}

Bước sinh câu trả lời nhận toàn bộ output của verification thay vì chỉ nhận
một danh sách rule. Answer intent được suy ra từ query text và
\texttt{query\_analysis}. Các claim direct hoặc counterfactual luôn được ưu
tiên ánh xạ thành \texttt{CLAIM\_VERIFICATION}, nhằm giữ nguyên semantics của
Step 4.

Đối với các câu hỏi còn lại, hệ thống nhận diện các mẫu ngôn ngữ liên quan
đến:

\begin{itemize}
    \item hệ quả cuối;
    \item nguyên nhân ban đầu;
    \item sự kiện cầu nối;
    \item giải thích chuỗi;
    \item câu hỏi có/không;
    \item phân nhánh;
    \item hội tụ.
\end{itemize}

Quá trình này không sử dụng \texttt{question\_type}. Do đó, cùng một cơ chế
có thể được áp dụng cho truy vấn ngoài benchmark nếu cách diễn đạt chứa tín
hiệu tương ứng.

\subsection{Lựa chọn path và evidence cho câu trả lời}

Step 5 lựa chọn path trước, sau đó lựa chọn evidence theo các path đã chọn.
Primary path luôn được giữ nếu tồn tại.

Đối với câu hỏi tuyến tính, evidence thuộc primary path được đưa lên trước
và giữ theo thứ tự causal. Đối với branch reasoning, hệ thống giữ các path
có chung prefix với primary path. Đối với convergence reasoning, hệ thống
giữ các path có chung suffix. Cơ chế này tránh trường hợp chỉ chọn một path
và làm mất các nhánh cần thiết để trả lời.

Evidence cuối cùng được lựa chọn từ:

\begin{enumerate}
    \item rule thuộc primary path;
    \item rule thuộc các path cấu trúc liên quan;
    \item evidence bổ sung đã được Step 4 xác minh;
    \item retrieval evidence fallback nếu verification không giữ được
    evidence nào.
\end{enumerate}

Các rule trùng nhau được loại bỏ theo rule ID. Khi tạo citation, các evidence
biểu diễn cùng một causal edge tiếp tục được khử trùng để tránh liệt kê lặp
lại nhiều căn cứ tương đương.

\subsection{Sinh câu trả lời extractive}

Mặc dù implementation hỗ trợ các provider bên ngoài như Ollama, OpenAI
compatible API và Gemini, thực nghiệm chính sử dụng extractive provider.
Lựa chọn này giúp kết quả có tính tất định và tách ảnh hưởng của retrieval
và verification khỏi biến thiên của mô hình sinh.

Extractive provider tạo câu trả lời theo answer intent:

\begin{description}
    \item[\texttt{FINAL\_OUTCOME}] Sử dụng nội dung \texttt{effect} đầy đủ
    của rule cuối trên primary path. Nếu không tìm được rule tương ứng, hệ
    thống fallback về outcome event name.

    \item[\texttt{INITIAL\_CAUSE}] Trả về event đầu tiên và chuỗi nhân quả
    đầy đủ.

    \item[\texttt{BRIDGE\_EVENT}] Trả về một hoặc nhiều mediator nằm giữa
    seed và outcome.

    \item[\texttt{CHAIN\_EXPLANATION}] Trình bày từng hop theo thứ tự, kèm
    citation của rule hỗ trợ hop đó.

    \item[\texttt{YES\_NO}] Mở đầu bằng ``Có'', ``Không'' hoặc ``Chưa đủ căn
    cứ'', phù hợp với final decision.

    \item[\texttt{BRANCH}] Tổng hợp outcome của các path có chung prefix.

    \item[\texttt{CONVERGENCE}] Tổng hợp các điều kiện đầu vào của những path
    có chung suffix.

    \item[\texttt{CLAIM\_VERIFICATION}] Trả lời trực tiếp claim và sử dụng
    decision explanation của Step 4.
\end{description}

Ví dụ, với primary path:

\[
\text{Cố ý thực hiện tội phạm nhưng không hoàn thành}
\rightarrow
\text{Phạm tội chưa đạt}
\rightarrow
\text{Chịu trách nhiệm hình sự},
\]

rule cuối chứa effect ``phải chịu trách nhiệm hình sự về tội phạm chưa đạt''.
Do đó, câu trả lời sử dụng effect đầy đủ này thay vì chỉ sử dụng event label
chung ``Chịu trách nhiệm hình sự''.

\subsection{Citation và kiểm tra đầu ra}

Mỗi evidence được gán một chỉ số \([E_i]\). Answer validator trích xuất các
citation xuất hiện trong câu trả lời và kiểm tra chúng với danh sách evidence
đã chọn.

Quá trình kiểm tra gồm:

\begin{itemize}
    \item loại citation không tồn tại;
    \item ưu tiên citation thuộc primary path;
    \item bổ sung citation nếu câu trả lời sử dụng evidence nhưng thiếu đánh
    dấu;
    \item lưu validation warning;
    \item ánh xạ evidence citation về article ID cho bước đánh giá.
\end{itemize}

Output cuối cùng lưu đồng thời citation dạng chỉ số evidence và căn cứ dạng
``Điều \(X\)''. Điều này cho phép kiểm tra cả grounding nội bộ của câu trả
lời và độ chính xác so với gold citation.

\subsection{Provenance và khả năng tái lập}

Mỗi artifact lưu phiên bản của bước tạo ra nó. Trong thí nghiệm chính:

\begin{itemize}
    \item Step 4 sử dụng phiên bản
    \texttt{4.1-query-aware-legal-scm};

    \item Step 5 sử dụng phiên bản
    \texttt{5.1-query-aware-grounded-answer};

    \item batch runner sử dụng phiên bản
    \texttt{2.1-query-aware-answer};

    \item evaluator sử dụng phiên bản
    \texttt{2.0-current-pipeline-schema}.
\end{itemize}

Generation metadata còn lưu:

\begin{itemize}
    \item answer intent;
    \item phiên bản intent detector;
    \item counterfactual signal source;
    \item Step 4 và Step 5 version;
    \item structural result hoặc structural fallback;
    \item provider và model;
    \item thời gian thực thi;
    \item số evidence và path được sử dụng;
    \item danh sách rule thuộc primary path.
\end{itemize}

Nhờ provenance này, kết quả từ các phiên bản pipeline khác nhau có thể được
phân biệt, tránh trường hợp so sánh artifact cũ và mới dưới cùng một nhãn
thực nghiệm.
